\documentclass[conference]{IEEEtran}
\usepackage{amsmath,amssymb}
\usepackage{graphicx}
\usepackage{booktabs}
\usepackage{makecell}
\usepackage{xcolor}
\usepackage{cite}
\usepackage{url}

% Figures are in the project Loop page; export them into figures/ and replace the boxes.
\newcommand{\figbox}[1]{\fbox{\parbox[c][3.0cm][c]{0.92\linewidth}{\centering\small\textit{#1}}}}
\newcommand{\todo}[1]{\textcolor{red}{[TODO: #1]}}

\begin{document}

\title{Narrow Frequency Range and Distorted Anti-spoofing}

\author{\IEEEauthorblockN{\todo{Authors}}
\IEEEauthorblockA{\todo{Affiliation}\\ \todo{Email}}}

\maketitle

%===============================================================
\begin{abstract}
Anti-spoofing systems are trained and benchmarked on wideband
speech, yet a large share of real calls reaches them through a
telephone channel that keeps only 300--3400\,Hz. We study how this narrow
frequency range (NFR) affects seven anti-spoofing systems: DF-Arena-1B, RawNet2, LFCC-LCNN,
AASIST, Nes2Net, TCM and SLS. NFR versions of the ASVspoof~2019 LA
evaluation set, In-the-Wild and a Famous Figures (Donald Trump) set are
created in two ways. The first uses five DSP band-limiting methods:
Parks--McClellan, Kaiser FIR, elliptic IIR, polyphase decimation and FFT
overlap-add. The second uses the ITU-T G.191 toolkit to apply IRS handset
shaping, P.56 level normalisation and G.711 $\mu$-law coding. In-domain,
band-limiting raises the EER of RawNet2 from 4.42\% to 24--26\% and of LCNN
from 6.54\% to about 18\%. SSL-based systems remain below 4\%. Out-of-domain,
the EER of Nes2Net, TCM and SLS on In-the-Wild rises from 10--14\% to
22--33\%, and at a fixed threshold 84--94\% of bona fide speech is
rejected. The choice of DSP method changes the EER by at most 3 percentage
points. A sub-band ablation shows that removing content below 300\,Hz,
rather than above 3.4\,kHz, causes most of the degradation.
\end{abstract}

\begin{IEEEkeywords}
anti-spoofing, deepfake detection, narrowband speech, telephony, ITU-T G.191, band-limiting
\end{IEEEkeywords}

%===============================================================
\section{Introduction}

Synthetic and converted speech is increasingly used in phone-based fraud,
where audio passes through a telephone channel before any anti-spoofing system sees
it. Such channels limit the signal to roughly 300--3400\,Hz at an 8\,kHz
sampling rate and add handset colouring, level control and codec
distortion. Anti-spoofing systems, however, are developed on
wideband 16\,kHz audio such as ASVspoof~2019~\cite{wang2020asvspoof}, and
their robustness to this narrow frequency range (NFR) is rarely measured
in a controlled way.

This paper measures how NFR conversion affects a broad set of anti-spoofing systems, from
raw-waveform and spectral baselines to models built on self-supervised
(SSL) speech representations. The contributions are:
\begin{itemize}
  \item \textbf{Controlled NFR data.} We convert speech to the telephone
  band in two ways: five DSP filtering methods, and the ITU-T G.191
  toolkit~\cite{itu_g191} used by operators to test call quality. Each
  condition changes a single factor relative to a control, so the cause of
  any change in performance can be identified.
  \item \textbf{Benchmark.} We evaluate seven anti-spoofing systems on one in-domain set
  (ASVspoof~2019 LA evaluation) and two out-of-domain sets (In-the-Wild
  and Famous Figures -- Donald Trump).
  \item \textbf{Analysis.} We break errors down by attack, by false
  positive and false rejection rate at a fixed threshold, by score
  distribution, and by which sub-band is removed.
\end{itemize}

%===============================================================
\section{Narrow Frequency Band Data Creation}

\subsection{NFR Methods}
\label{sec:nfr}

We implement five methods that keep the 300\,Hz--3.4\,kHz band, chosen to
cover the main families of practical band-limiting.

\textbf{F1 -- Parks--McClellan (equiripple FIR)}~\cite{mcclellan1973}
chooses the filter coefficients that minimise the maximum error over the
spectrum. The residual error is spread evenly rather than concentrated at
a few frequencies, giving the sharpest cut for a given filter length.

\textbf{F2 -- Kaiser window FIR}~\cite{kaiser1974} truncates the ideal
infinite-length filter and tapers it with a Kaiser window to suppress
ringing. It needs a somewhat longer filter than F1 for the same quality,
but its closed-form design cannot fail, which suits automatically
generated filters.

\textbf{F3 -- Elliptic IIR} feeds past outputs back into the filter and is
therefore about eight times cheaper than the FIR designs. This is why it
is common in real-time processing on phones and embedded hardware. Its
non-linear phase delays frequencies by different amounts, which smears
timing near the band edges.

\textbf{F4 -- Polyphase decimation} lowers the sampling rate instead of
filtering, so content above the new Nyquist frequency can no longer be
represented. It is the fastest method and also reduces file size. If the
anti-aliasing stage is weak, however, high-frequency energy folds back
into the kept band, where it cannot be removed.

\textbf{F5 -- FFT overlap-add} applies the same filter as F2 block-wise in
the frequency domain. The output is the same as F2, but it is about ten
times faster for long filters.

\textbf{F0 -- control} is the source signal scaled by the same constant
$-6$\,dB headroom applied in F1--F5, so F0 differs from F1--F5 only in the
filtering.

Fig.~\ref{fig:freqresp} shows the magnitude response of each method and
Table~\ref{tab:filters} compares them on the same speech. Kaiser FIR and
polyphase decimation give the strongest stop-band rejection above 3.4\,kHz
(161\,dB) and the flattest passband ($\le$0.02\,dB). The elliptic IIR has
the narrowest transition (45\,Hz) and is the fastest (11\,ms per 6\,s),
but it has the weakest rejection (65\,dB above the band, 61\,dB below) and
1\,dB of passband ripple.

\begin{table}[t]
\centering
\caption{The five NFR methods when keeping 300\,Hz--3.4\,kHz of the same utterance. Transition width is measured from $-3$ to $-60$\,dB (both edges averaged); time is for 6\,s of audio including filter design. Best values in bold.}
\label{tab:filters}
\setlength{\tabcolsep}{3pt}
\begin{tabular}{lccccc}
\toprule
Method & \makecell{Rejection\\$>$3.4\,kHz} & \makecell{Rejection\\$<$300\,Hz} & \makecell{Passband\\ripple} & \makecell{Trans.\\width} & Time \\
\midrule
F1 Parks--McClellan & 122\,dB & \textbf{86\,dB} & 0.39\,dB & 74\,Hz & 199\,ms \\
F2 Kaiser FIR       & \textbf{161\,dB} & \textbf{87\,dB} & \textbf{0.02\,dB} & 64\,Hz & 241\,ms \\
F3 Elliptic IIR     & 65\,dB & 61\,dB & 1.00\,dB & \textbf{45\,Hz} & \textbf{11\,ms} \\
F4 Polyphase dec.   & \textbf{161\,dB} & \textbf{86\,dB} & \textbf{0.01\,dB} & 64\,Hz & 38\,ms \\
F5 FFT overlap-add  & 86\,dB & 74\,dB & \textbf{0.01\,dB} & 65\,Hz & \textbf{11\,ms} \\
\bottomrule
\end{tabular}
\end{table}

\begin{figure}[t]
\centering
\figbox{Magnitude response of F1--F5; grey band = 300\,Hz--3.4\,kHz (Loop: ``What each method keeps and removes'').}
\caption{Magnitude response of the five NFR methods. The further a curve lies below 0\,dB outside the grey band, the more completely that method removes out-of-band content.}
\label{fig:freqresp}
\end{figure}

\textbf{Spectral inspection.} We inspect the outputs with constant-Q
log spectrograms~\cite{brown1991cqt}. A fixed STFT window of 2048 samples
(43\,ms at 48\,kHz, hop 256) resolves only about 23\,Hz at every frequency.
That is too coarse for the speaker's fundamental frequency of about
100--140\,Hz, which therefore appears blurred. We instead use $Q=18$, 48
bins per octave, a 50\,Hz--20\,kHz range and a 5\,ms hop. The window then
lasts about 150\,ms at 120\,Hz and 4.5\,ms at 4\,kHz, so F0 and its
harmonics appear as distinct lines while high-frequency timing stays
sharp (Fig.~\ref{fig:cqt}).

\begin{figure}[t]
\centering
\figbox{Constant-Q log spectrograms of p232\_003.wav (48\,kHz, 7.2\,s): original, F1--F5, and change from the original per method.}
\caption{Constant-Q log spectrograms of one utterance before and after each NFR method.}
\label{fig:cqt}
\end{figure}

\subsection{Telephony Simulation with ITU-T G.191}
\label{sec:itu}

To reproduce a real telephone path we use the ITU-T G.191 Software Tool
Library (STL)~\cite{itu_g191}. Telephone operators use these tools to test
call quality, so the resulting conditions follow an industry standard. The
library is free and deterministic, and it runs on our compute cluster. We
define three nested conditions, each adding one step:
\begin{itemize}
  \item \textbf{P0} -- control: active speech level normalised to
  $-26$\,dBov with ITU-T P.56~\cite{itu_p56}, 16\,kHz;
  \item \textbf{B1} -- narrowband telephone: resampled to 8\,kHz and
  shaped to 300--3400\,Hz by the IRS handset response
  (ITU-T P.48~\cite{itu_p48});
  \item \textbf{C1} -- B1 followed by G.711 $\mu$-law
  coding~\cite{itu_g711}, matching a landline call.
\end{itemize}
\textbf{RAW} denotes the unprocessed audio.

\textbf{Physical recordings.} To compare the simulation with real
capture, one utterance was recorded over a live phone call and on two
separate devices (A and B). The recordings were time-aligned to Device~A
and compared with constant-Q spectrograms (Fig.~\ref{fig:phys}).

\begin{figure}[t]
\centering
\figbox{Constant-Q spectrograms (48 bins/oct, 16\,kHz) of Utterance~3: phone call, Device~A, Device~B.}
\caption{The same utterance captured over a phone call and on two devices.}
\label{fig:phys}
\end{figure}

%===============================================================
\section{Anti-spoofing Systems}

\subsection{Architecture}

We evaluate seven publicly available anti-spoofing systems. All except DF-Arena-1B were
trained on the ASVspoof~2019 LA training partition, so its evaluation
partition is in-domain for them. Each model is run with its released
checkpoint and its own input length. Shorter clips are tile-repeated to
fill that window.

\textbf{RawNet2}~\cite{tak2021rawnet2} is the official ASVspoof~2021
baseline~\cite{yamagishi2021asvspoof}. It operates directly on the raw
waveform with a learnable sinc-filter front end followed by residual
blocks and a GRU (about 25M parameters, trained from scratch, 4\,s input).

\textbf{LFCC-LCNN}~\cite{lavrentyeva2019stc,yamagishi2021asvspoof} is the
second ASVspoof~2021 baseline. It computes linear-frequency cepstral
coefficients inside the model and classifies them with a light CNN using
max-feature-map activations.

\textbf{AASIST}~\cite{jung2022aasist,tak2022ssl} uses an XLS-R
(300M)~\cite{babu2022xlsr} SSL front end with the AASIST spectro-temporal
graph-attention back end (0.45M parameters).

\textbf{TCM}~\cite{truong2024tcm} feeds XLS-R features to a Conformer whose
multi-head self-attention models temporal and channel dependencies jointly
(4 encoders, 4 heads, embedding size 144).

\textbf{SLS}~\cite{zhang2024sls} learns to select and weight the most
spoof-sensitive XLS-R layers before classification. Its 23.4M-parameter
back end is the largest in the suite.

\textbf{Nes2Net-X}~\cite{liu2025nes2net} uses a nested Res2Net back end
directly on the SSL features without a dimensionality-reduction layer
(0.51M parameters).

\textbf{DF-Arena-1B}~\cite{dfarena} is a publicly released
1B-parameter anti-spoofing system from the Speech Arena project.
\todo{architecture details}

%===============================================================
\section{Experimental Setup}

\subsection{Pipeline}

\textbf{DSP route.} Each source file is processed by F0 and F1--F5
(Section~\ref{sec:nfr}) and kept at 16\,kHz.

\textbf{ITU route.} Real and fake audio are converted to mono and resampled
to 16\,kHz (P0). For B1 and C1 they are further downsampled to 8\,kHz,
shaped by the IRS8 handset filter, level-normalised with P.56 and, for C1,
coded with G.711. Fig.~\ref{fig:itu} shows a real and a fake clip before
and after this pipeline.

\begin{figure}[t]
\centering
\figbox{Constant-Q spectrograms before/after the ITU pipeline: real Donald\_Trump\_00386 (16\,kHz $\rightarrow$ 8\,kHz) and fake CosyVoice2 Donald\_Trump\_00353 (24\,kHz $\rightarrow$ 8\,kHz).}
\caption{A real and a fake clip before and after the ITU-T telephone pipeline.}
\label{fig:itu}
\end{figure}

\textbf{Sub-band ablation.} To determine which part of the removed
spectrum matters, we remove one side at a time. HP100, HP300, HP500 and
HP1000 high-pass at 100, 300, 500 and 1000\,Hz, and LP3400 low-passes at
3.4\,kHz. These are compared with F3, which removes both sides.

\textbf{Metrics.} We report the equal error rate (EER) with a 95\%
confidence interval, AUC and $d'$. To measure how a deployed operating
point behaves, we fix each model's threshold at its EER point on RAW
ASVspoof~2019 LA. We then report the false positive rate (FPR, spoof
accepted) and false rejection rate (FRR, bona fide rejected) under every
condition, together with the means and spreads of the real and fake
scores.

\subsection{In-domain Dataset}

\textbf{ASVspoof~2019 LA evaluation}~\cite{wang2020asvspoof}: 71{,}237
trials covering attacks A07--A19. All seven systems, apart from DF-Arena-1B,
were trained on its training partition. Every condition uses the same
trial list and labels.

\subsection{Out-of-domain Datasets}

\textbf{In-the-Wild (ITW)}~\cite{muller2022itw}: real and deepfake
recordings of public figures collected from online media, with
uncontrolled recording conditions.

\textbf{Famous Figures -- Donald Trump (FF)} \todo{cite}: bona fide and
synthetic speech of a single speaker produced by several modern TTS and
voice-cloning systems.

%===============================================================
\section{Results and Analysis}

\subsection{DF-Arena-1B on ASVspoof~2019 LA}

\begin{table}[t]
\centering
\caption{DF-Arena-1B on ASVspoof~2019 LA eval under the ITU conditions.}
\label{tab:dfa_itu}
\begin{tabular}{lccc}
\toprule
 & P0 & B1 & C1 \\
\midrule
EER (\%) & \textbf{1.455} & 2.243 & 2.325 \\
95\% CI  & [1.36, 1.55] & [2.12, 2.37] & [2.22, 2.49] \\
AUC      & 0.99743 & 0.99412 & 0.99374 \\
$d'$     & 5.185 & 4.531 & 4.512 \\
\bottomrule
\end{tabular}
\end{table}

\begin{table}[t]
\centering
\caption{DF-Arena-1B per-attack EER (\%), ASVspoof~2019 LA eval.}
\label{tab:dfa_attack}
\begin{tabular}{lcccl}
\toprule
Attack & P0 & B1 & C1 & \\
\midrule
A10 & 8.101 & 12.675 & 13.310 & hardest throughout \\
A16 & 1.076 & 2.567 & 2.788 & $2.4\times$ \\
A12 & 0.367 & 0.856 & 0.896 & $2.3\times$ \\
A19 & 0.530 & 0.750 & 0.774 & \\
A18 & 0.465 & 0.879 & 0.937 & \\
A13 & 1.165 & 1.117 & 1.222 & unchanged \\
A08, A09 & 0.000 & 0.000 & 0.000 & unaffected \\
\bottomrule
\end{tabular}
\end{table}

\begin{table}[t]
\centering
\caption{DF-Arena-1B on ASVspoof~2019 LA eval under the five DSP methods ($n=71{,}237$).}
\label{tab:dfa_dsp}
\begin{tabular}{lcccc}
\toprule
Cond. & EER (\%) & 95\% CI & AUC & $d'$ \\
\midrule
F1 Parks--McClellan & 4.870 & [4.69, 5.09] & 0.9854 & 3.730 \\
F2 Kaiser FIR       & 4.079 & [3.92, 4.24] & 0.9890 & 3.956 \\
F3 Elliptic IIR     & \textbf{3.656} & [3.51, 3.84] & 0.9900 & 3.985 \\
F4 Polyphase        & 3.805 & [3.64, 3.97] & 0.9897 & 3.970 \\
F5 FFT overlap-add  & 3.970 & [3.81, 4.13] & 0.9894 & 3.986 \\
\bottomrule
\end{tabular}
\end{table}

Telephone shaping raises the EER of DF-Arena-1B from 1.455\% (P0) to
2.243\% (B1), and $d'$ falls from 5.19 to 4.53
(Table~\ref{tab:dfa_itu}). Adding G.711 changes little (2.325\%, with
overlapping confidence intervals), so the codec is not the main source of
degradation. Per attack (Table~\ref{tab:dfa_attack}), A10 dominates the
error in every condition. A16 and A12 more than double, A13 is unaffected
and A08/A09 remain perfectly detected. Under the DSP methods the EER is
3.66--4.87\% (Table~\ref{tab:dfa_dsp}). F3 gives the lowest EER and F1 the
highest.

\subsection{Six Countermeasures across Datasets}

\begin{table*}[t]
\centering
\caption{EER (\%) of six anti-spoofing systems on the in-domain set (ASVspoof~2019 LA eval) and the two out-of-domain sets. RN2: RawNet2; N2N: Nes2Net.}
\label{tab:eer}
\setlength{\tabcolsep}{3.1pt}
\begin{tabular}{ll|cccccc|cccccc|cccccc}
\toprule
 & & \multicolumn{6}{c|}{ASVspoof~2019 LA} & \multicolumn{6}{c|}{In-the-Wild} & \multicolumn{6}{c}{Famous Figures -- Trump} \\
Cond. & What & RN2 & LCNN & AASIST & N2N & TCM & SLS & RN2 & LCNN & AASIST & N2N & TCM & SLS & RN2 & LCNN & AASIST & N2N & TCM & SLS \\
\midrule
RAW & original     & 4.60 & 6.35 & 0.23 & 0.14 & 0.15 & 0.23 & 48.99 & 29.48 & 11.22 & 8.41 & 7.94 & 7.45 & 46.07 & 28.21 & 58.58 & 43.04 & 39.76 & 40.28 \\
P0  & level only   & 6.14 & 6.53 & 0.20 & 0.15 & 0.23 & 0.39 & 39.41 & 27.70 & 10.39 & 9.14 & 7.96 & 6.78 & 49.41 & 28.98 & 59.15 & 39.88 & 36.35 & 36.64 \\
B1  & phone filter & 35.73 & 18.14 & 0.24 & 0.78 & 3.15 & 3.13 & 41.41 & 21.96 & 13.12 & 19.49 & 22.74 & 16.18 & 38.34 & 39.00 & 55.78 & 47.62 & 44.53 & 45.21 \\
C1  & + codec      & 35.69 & 53.65 & 0.25 & 0.77 & 3.18 & 3.11 & 41.36 & 28.23 & 13.07 & 19.58 & 22.96 & 16.38 & 38.43 & 33.40 & 55.89 & 47.86 & 44.86 & 45.18 \\
\midrule
F0  & control      & 4.42 & 6.54 & 0.20 & 0.12 & 0.22 & 0.32 & 51.82 & 29.42 & 13.83 & 13.96 & 11.34 & 10.16 & 44.89 & 28.37 & 58.41 & 42.72 & 42.81 & 47.28 \\
F1  & Parks--McC.  & 24.79 & 18.25 & 0.32 & 1.21 & 3.22 & 3.53 & 50.61 & 29.99 & 16.35 & 22.08 & 30.21 & 23.30 & 28.63 & 42.99 & 56.06 & 50.53 & 47.73 & 49.58 \\
F2  & Kaiser FIR   & 24.46 & 18.21 & 0.29 & 0.97 & 2.81 & 3.29 & 50.76 & 29.83 & 16.25 & 21.79 & 30.04 & 23.09 & 28.61 & 42.87 & 56.21 & 50.76 & 47.79 & 49.58 \\
F3  & Elliptic IIR & 26.35 & 18.34 & 0.29 & 0.86 & 2.56 & 3.05 & 50.54 & 30.39 & 17.10 & 24.11 & 33.08 & 25.49 & 28.40 & 41.90 & 56.14 & 50.82 & 46.91 & 48.27 \\
F4  & Polyphase    & 25.44 & 18.45 & 0.31 & 0.94 & 2.64 & 3.18 & 50.52 & 29.99 & 16.90 & 23.48 & 32.70 & 25.18 & 28.50 & 41.37 & 56.04 & 51.09 & 47.22 & 48.99 \\
F5  & FFT overlap  & 24.16 & 18.27 & 0.30 & 0.94 & 2.75 & 3.21 & 51.21 & 29.85 & 16.16 & 21.80 & 30.01 & 22.96 & 28.48 & 43.31 & 56.06 & 50.26 & 47.45 & 49.73 \\
\bottomrule
\end{tabular}
\end{table*}

\textbf{In-domain.} Table~\ref{tab:eer} shows two distinct groups. The
models without an SSL front end degrade sharply. RawNet2 rises from 4.42\%
(F0) to 24.16--26.35\% under F1--F5 and to 35.7\% under B1/C1. LCNN rises
from 6.54\% to about 18.3\%, and it is the only model sensitive to the
codec: its EER rises from 18.14\% at B1 to 53.65\% at C1. The SSL-based
models stay usable. AASIST barely moves (at most 0.32\%), Nes2Net reaches
0.86--1.21\%, and TCM and SLS rise roughly tenfold to 2.56--3.53\%. Per-file
level normalisation alone (RAW$\rightarrow$P0) costs RawNet2 1.5
percentage points, whereas the constant gain of F0 is nearly free (4.60\%
$\rightarrow$ 4.42\%). We therefore compare each ITU condition with P0 and
each DSP method with F0.

\textbf{Out-of-domain.} On In-the-Wild, the models with the lowest
in-domain EERs degrade the most. Relative to F0, under F1--F5 Nes2Net goes
from 13.96\% to 21.79--24.11\%, TCM from 11.34\% to 30.01--33.08\%, SLS from
10.16\% to 22.96--25.49\% and AASIST from 13.83\% to 16.16--17.10\%. The ITU
path degrades TCM and SLS less (22.7\% and 16.2\% at B1) than the DSP
filters. RawNet2 is close to chance in every condition. On Famous Figures,
every model is already near chance before processing (28--59\% EER), so
NFR effects cannot be measured reliably on this set.

\textbf{Choice of NFR method.} Within a model and dataset, the five DSP
methods give similar EERs. The largest spread is 2.19 percentage points
(RawNet2, ASVspoof) and 3.07 percentage points (TCM, In-the-Wild), small
compared with the effect of band-limiting itself. The ranking is not
consistent across models: F3 is the worst method for RawNet2 but the best
for TCM and SLS on ASVspoof.

\subsection{Error Analysis at a Fixed Threshold}

\begin{table*}[t]
\centering
\caption{FPR / FRR (\%) with each model's threshold fixed at its EER point on RAW ASVspoof~2019 LA.}
\label{tab:fprfrr}
\setlength{\tabcolsep}{2.5pt}
\begin{tabular}{ll|cccccc|cccccc}
\toprule
 & & \multicolumn{6}{c|}{ASVspoof~2019 LA} & \multicolumn{6}{c}{In-the-Wild} \\
Cond. & What & RawNet2 & LCNN & AASIST & Nes2Net & TCM & SLS & RawNet2 & LCNN & AASIST & Nes2Net & TCM & SLS \\
\midrule
RAW & original     & 4.60/4.60 & 6.35/6.35 & 0.23/0.23 & 0.14/0.14 & 0.15/0.15 & 0.23/0.23 & 5.92/88.21 & 9.13/64.43 & 0.29/77.69 & 0.17/64.88 & 0.72/53.04 & 0.30/60.54 \\
P0  & level only   & 3.17/12.43 & 9.43/2.94 & 0.28/0.16 & 0.12/0.23 & 0.19/0.30 & 0.20/0.76 & 2.49/88.62 & 8.99/63.95 & 0.27/74.08 & 0.10/66.24 & 0.64/49.92 & 0.25/59.06 \\
B1  & phone filter & 34.96/36.66 & 22.87/1.82 & 0.14/0.58 & 0.10/16.41 & 0.98/14.33 & 0.34/39.76 & 17.62/65.47 & 13.07/40.50 & 0.32/82.40 & 0.15/94.89 & 1.67/80.56 & 0.43/90.39 \\
C1  & + codec      & 34.93/36.46 & 12.21/76.61 & 0.14/0.61 & 0.09/16.67 & 0.99/14.41 & 0.34/39.69 & 17.69/65.44 & 26.03/32.37 & 0.32/82.51 & 0.15/94.97 & 1.68/80.71 & 0.42/90.45 \\
\midrule
F0  & control      & 2.78/8.27 & 8.44/3.81 & 0.26/0.15 & 0.12/0.14 & 0.19/0.24 & 0.22/0.52 & 18.77/80.15 & 13.64/56.27 & 0.29/82.12 & 0.27/81.00 & 1.11/65.59 & 0.36/75.07 \\
F1  & Parks--McC.  & 3.61/60.31 & 18.22/18.29 & 0.09/4.68 & 0.09/15.47 & 0.35/18.22 & 0.36/35.23 & 5.21/88.53 & 7.41/84.33 & 0.25/93.44 & 0.56/94.00 & 5.39/84.51 & 2.24/92.14 \\
F2  & Kaiser FIR   & 3.60/59.90 & 18.86/16.07 & 0.10/2.56 & 0.10/10.81 & 0.35/15.76 & 0.36/32.70 & 5.23/88.53 & 7.33/84.40 & 0.25/92.80 & 0.54/93.68 & 5.32/84.29 & 2.23/91.69 \\
F3  & Elliptic IIR & 4.28/61.36 & 17.52/20.86 & 0.09/2.39 & 0.11/10.89 & 0.51/11.28 & 0.39/31.38 & 4.97/88.25 & 7.69/84.73 & 0.27/93.52 & 0.66/93.44 & 7.11/83.60 & 2.76/91.61 \\
F4  & Polyphase    & 5.17/56.64 & 17.57/21.36 & 0.10/2.68 & 0.10/11.54 & 0.48/11.57 & 0.39/32.22 & 5.69/87.44 & 7.26/85.98 & 0.25/93.73 & 0.63/94.24 & 7.00/83.82 & 2.67/92.28 \\
F5  & FFT overlap  & 3.32/60.72 & 19.11/15.45 & 0.10/2.58 & 0.10/10.78 & 0.36/15.38 & 0.36/33.11 & 5.13/88.98 & 7.82/83.18 & 0.25/92.87 & 0.60/93.61 & 5.12/84.43 & 2.22/91.79 \\
\bottomrule
\end{tabular}
\end{table*}

EER alone hides how a deployed anti-spoofing system would fail. With the threshold
fixed (Table~\ref{tab:fprfrr}), NFR conversion shows up almost entirely as
bona fide speech being rejected, while spoofs are still caught. On
ASVspoof, RawNet2 rejects 57--61\% of real speech under F1--F5 while its
FPR stays at 3--5\%. SLS rejects 40\% of real speech under B1 at an FPR of
0.34\%. On In-the-Wild the threshold transfers poorly even without
processing: AASIST, Nes2Net, TCM and SLS already reject 53--78\% of real
clips on RAW, and under F1--F5 this rises to 84--94\% with FPR below
7.2\%.

\begin{table}[t]
\centering
\caption{In-the-Wild score analysis. Top: real-clip scores against the fixed threshold. Bottom: separation of real and fake. Sep.\ = gap / pooled spread.}
\label{tab:scores}
\setlength{\tabcolsep}{2.6pt}
\begin{tabular}{l|ccc|ccc}
\toprule
 & \multicolumn{3}{c|}{RAW (original)} & \multicolumn{3}{c}{F5 (300--3400\,Hz)} \\
Model & \makecell{Real $\mu$ / thr.} & \makecell{Below thr.\\($\sigma$)} & FRR & \makecell{Real $\mu$ / thr.} & \makecell{Below thr.\\($\sigma$)} & FRR \\
\midrule
AASIST  & $-$0.06 / +2.00 & 0.77 & 77.69 & $-$1.83 / +2.00 & 1.48 & 92.87 \\
Nes2Net & +1.56 / +3.13 & 0.50 & 64.88 & $-$1.17 / +3.13 & 1.89 & 93.61 \\
TCM     & +0.09 / +0.43 & 0.16 & 53.04 & $-$1.40 / +0.43 & 1.05 & 84.43 \\
SLS     & $-$3.98 / $-$0.86 & 0.73 & 60.54 & $-$7.03 / $-$0.86 & 1.75 & 91.79 \\
\midrule
Model & Gap & Sep. & EER & Gap & Sep. & EER \\
\midrule
AASIST  & 5.3  & 2.43 & 11.22 & 3.6 & 1.74 & 16.16 \\
Nes2Net & 5.0  & 2.20 & 8.41  & 2.0 & 1.13 & 21.80 \\
TCM     & 4.1  & 2.55 & 7.94  & 1.4 & 0.84 & 30.01 \\
SLS     & 10.3 & 3.09 & 7.45  & 4.1 & 1.25 & 22.96 \\
\bottomrule
\end{tabular}
\end{table}

Table~\ref{tab:scores} explains this behaviour. On RAW In-the-Wild audio,
the mean real-clip score of the SSL models already sits 0.16--0.77
standard deviations below the ASVspoof threshold. Narrowband conversion
has two effects. First, real-clip scores move further towards the spoof
side, to 1.05--1.89 standard deviations below the threshold. Second, the
real and fake distributions move closer together: separation falls from
2.43 to 1.74 (AASIST), 2.20 to 1.13 (Nes2Net), 2.55 to 0.84 (TCM) and 3.09
to 1.25 (SLS). The first effect could be partly offset by recalibrating
the threshold. The second cannot, and it is what drives the EER increase.

\subsection{Sub-band Ablation}

\begin{table}[t]
\centering
\caption{EER (\%) when one side of the spectrum is removed. F3 removes both sides.}
\label{tab:subband}
\setlength{\tabcolsep}{2.8pt}
\begin{tabular}{llcccccc}
\toprule
Cond. & Removed & RN2 & LCNN & AASIST & N2N & TCM & SLS \\
\midrule
\multicolumn{8}{l}{\textit{ASVspoof~2019 LA}} \\
raw    & nothing     & 4.42 & 6.54 & 0.20 & 0.12 & 0.22 & 0.32 \\
HP100  & $<$100\,Hz  & 4.93 & 6.68 & 0.20 & 0.15 & 0.24 & 0.38 \\
HP300  & $<$300\,Hz  & 18.64 & 15.00 & 0.26 & 0.31 & 1.47 & 2.11 \\
HP500  & $<$500\,Hz  & 26.78 & 18.15 & 0.44 & 0.62 & 5.30 & 3.29 \\
HP1000 & $<$1\,kHz   & 38.99 & 14.88 & 0.91 & 1.41 & 11.87 & 8.28 \\
LP3400 & $>$3.4\,kHz & 5.03 & 13.72 & 0.16 & 0.20 & 0.41 & 0.61 \\
F3     & both        & 26.35 & 18.34 & 0.29 & 0.86 & 2.56 & 3.05 \\
\midrule
\multicolumn{8}{l}{\textit{In-the-Wild}} \\
raw    & nothing     & 51.82 & 29.42 & 13.83 & 13.96 & 11.34 & 10.16 \\
HP100  & $<$100\,Hz  & 50.75 & 29.08 & 13.88 & 14.60 & 12.04 & 10.49 \\
HP300  & $<$300\,Hz  & 44.83 & 30.95 & 13.84 & 18.47 & 21.65 & 17.39 \\
HP500  & $<$500\,Hz  & 42.56 & 24.66 & 16.91 & 23.49 & 35.45 & 31.81 \\
HP1000 & $<$1\,kHz   & 39.94 & 35.07 & 23.25 & 28.98 & 42.78 & 43.10 \\
LP3400 & $>$3.4\,kHz & 58.24 & 30.58 & 16.52 & 18.15 & 16.24 & 13.19 \\
F3     & both        & 50.54 & 30.39 & 17.10 & 24.11 & 33.08 & 25.49 \\
\bottomrule
\end{tabular}
\end{table}

Table~\ref{tab:subband} separates the two halves of the telephone band
limit. On ASVspoof, removing everything above 3.4\,kHz (LP3400) has little
effect on any SSL model (at most 0.61\%) or on RawNet2 (5.03\%). Removing
the low end has a much larger effect, which grows steadily with the cutoff.
For RawNet2 the EER rises from 4.93\% (HP100) to 18.64\% (HP300) and 38.99\%
(HP1000). For TCM it rises from 0.24\% to 1.47\% and 11.87\%, and for SLS
from 0.38\% to 2.11\% and 8.28\%. HP300 alone accounts for most of the F3
degradation of RawNet2 (18.64\% vs.\ 26.35\%), TCM (1.47\% vs.\ 2.56\%) and
SLS (2.11\% vs.\ 3.05\%). LCNN is the exception: it is hurt by both sides
(15.00\% at HP300, 13.72\% at LP3400), consistent with its linear-frequency
cepstral front end weighting the high band.

The same pattern holds out-of-domain for the SSL models. On In-the-Wild,
TCM and SLS rise from 11.34\% and 10.16\% (raw) to 21.65\% and 17.39\% at
HP300 and to 42.78\% and 43.10\% at HP1000. At LP3400 they reach only 16.24\%
and 13.19\%. For both TCM and SLS the HP500 condition (35.45\%, 31.81\%) is
already worse than F3 itself (33.08\%, 25.49\%). RawNet2 is the outlier: its
near-chance EER on In-the-Wild falls as the high-pass cutoff rises (51.82\%
$\rightarrow$ 39.94\%) and worsens at LP3400 (58.24\%). Across models, the
content between roughly 100\,Hz and 1\,kHz, which holds the fundamental
frequency and the first harmonics, carries most of the evidence these anti-spoofing systems
rely on. This is exactly the content the telephone band removes.

%===============================================================
\section{Conclusion}

We built controlled narrowband versions of three anti-spoofing datasets,
using five DSP band-limiting methods and the ITU-T G.191 telephone chain,
and evaluated seven anti-spoofing systems on them. Narrow-band conversion
degrades every model. Raw-waveform and cepstral baselines fail even
in-domain, and SSL-based models, which are robust in-domain, lose most of
their margin out-of-domain, where nearly all bona fide speech is rejected
at a fixed operating point. The choice of filtering algorithm matters far
less than the band limit itself. Within that limit, losing the band below
300\,Hz is the main cause of degradation, not losing the band above
3.4\,kHz. These results suggest that training with telephone-band
augmentation and calibrating thresholds per channel are necessary before
deploying anti-spoofing systems on telephony audio. \todo{future work}

\bibliographystyle{IEEEtran}
\bibliography{refs}

\end{document}
